\chapter{Streams, Pipes and Redirection}\label{ch:sh-pipes}

\autoref{ch:sh-text} gave you a box of tools that each do one narrow thing.
This chapter is the plumbing that turns them into a language --- and it is the
part of Unix that has been copied by every system built since.

\section{Standard Input, Output and Error}\label{sec:sh-streams}

\begin{definition}[Standard streams]
	Every process starts with three open file descriptors. \vocab{Standard input}
	(fd~0) is where it reads from, \vocab{standard output} (fd~1) is where its
	results go, and \vocab{standard error} (fd~2) is where its complaints go.
	By default all three are your terminal.
\end{definition}

\begin{intuition}
	A program does not know or care what is on the other end of those three
	descriptors. It might be your keyboard, a file, or another program --- and the
	program is written identically either way. That indifference is the entire
	mechanism behind everything in this chapter.
\end{intuition}

Splitting output from errors looks fussy until the first time it saves you:

\begin{shell}
$ ls real.txt missing.txt
ls: cannot access 'missing.txt': No such file or directory     # this went to fd 2
real.txt                                                       # this went to fd 1
$ ls real.txt missing.txt 2>/dev/null
real.txt
$ ls real.txt missing.txt > out.txt
ls: cannot access 'missing.txt': No such file or directory     # errors still appear
\end{shell}

\begin{note}
	The second case is why error messages appear on screen even when you have
	redirected output to a file: they were never on fd~1. This is a feature ---
	a pipeline's data stays clean while its diagnostics still reach you.
\end{note}

\section{Redirection}\label{sec:sh-redirect}

Redirection attaches a stream to a file instead of the terminal. It is handled
by the shell before the program starts, so the program never sees the
\texttt{>}.

\begin{keys}[0.24]
	\texttt{> file}        & Send stdout to \file{file}, \textbf{destroying} its contents \\
	\texttt{>{}> file}       & Append to \file{file} instead                                \\
	\texttt{< file}        & Take stdin from \file{file}                                  \\
	\texttt{2> file}       & Send stderr to \file{file}                                   \\
	\texttt{2>{}> file}      & Append stderr                                                \\
	\texttt{2>\&1}         & Send stderr wherever stdout is currently going               \\
	\texttt{\&> file}      & Both streams to one file (Bash and Zsh)                      \\
	\texttt{> /dev/null}   & Discard --- \file{/dev/null} accepts anything and keeps none \\
	\texttt{n> file}       & Redirect an arbitrary descriptor                             \\
\end{keys}

\begin{note}[Order matters for \texttt{2>\&1}]
	\texttt{cmd > out.txt 2>\&1} sends both streams to the file. \texttt{cmd 2>\&1
		> out.txt} does \emph{not}: redirections are applied left to right, so stderr
	is pointed at the terminal (where stdout is at that moment) and only then is
	stdout moved to the file. The second form is almost always a bug, and it is
	subtle enough to survive code review.
\end{note}

\begin{moral}
	\texttt{>} truncates before the command even runs. So \cmd{sort file > file}
	empties the file and then sorts nothing, and there is no recovery. Write to a
	new name, or use a tool with an in-place option such as \cmd{sed -i}.
\end{moral}

\section{Pipes and Filters}\label{sec:sh-pipes}

\begin{definition}[Pipe]
	\texttt{a | b} connects \texttt{a}'s standard output directly to \texttt{b}'s
	standard input. Both run \emph{at the same time}, with the kernel buffering
	between them --- no temporary file is created, and \texttt{b} can start
	producing output before \texttt{a} has finished.
\end{definition}

\begin{definition}[Filter]
	A program that reads standard input and writes standard output, so that it can
	sit in the middle of a pipeline. \cmd{grep}, \cmd{sort}, \cmd{sed}, \cmd{tr},
	\cmd{uniq}, \cmd{wc} and \cmd{awk} are all filters, which is why they combine.
\end{definition}

\begin{eg}
	Reading a pipeline is a matter of reading it left to right, one arrow at a
	time:
\begin{shell}
$ ps aux | grep nvim | awk '{print $2}' | xargs kill
\end{shell}
	List every process; keep the lines mentioning \cmd{nvim}; take the second
	column of each, which is the process ID; hand those IDs to \cmd{kill}. Four
	single-purpose programs, no temporary files, no program aware of the others.
\end{eg}

\begin{note}
	The exit status of a pipeline is the status of its \emph{last} command, so a
	failure in the middle is silently ignored. \cmd{set -o pipefail}
	(\autoref{sec:sh-robust}) changes this, and any script with a pipeline in it
	should set it.
\end{note}

\begin{tryit}
	Build this up one stage at a time, pressing \key{<CR>} after each addition, and
	watch the output narrow:
	\cmd{cut -d: -f7 /etc/passwd}, then \cmd{| sort}, then \cmd{| uniq -c}, then
	\cmd{| sort -rn}. You have just answered ``which login shells are in use here,
	and how popular is each''.
\end{tryit}

\section{\texorpdfstring{\cmd{tee} and \cmd{xargs}}{tee and xargs}}\label{sec:sh-teexargs}

Two commands exist to patch the two places where pipes fall short.

\cmd{tee} solves ``I want this both on screen and in a file'':

\begin{keys}[0.26]
	\cmd{cmd | tee log.txt}       & Write to the file \emph{and} pass through       \\
	\cmd{cmd | tee -a log.txt}    & Append rather than overwrite                    \\
	\cmd{cmd | tee /dev/tty | c2} & Show the intermediate stage of a pipeline       \\
	\cmd{echo x | sudo tee /etc/f} & Write to a root-owned file --- \cmd{sudo cmd > file} does not work, because the shell opens the file, not \cmd{sudo} \\
\end{keys}

\cmd{xargs} solves the opposite problem: many programs take filenames as
\emph{arguments}, not on standard input.

\begin{keys}[0.30]
	\cmd{... | xargs rm}             & Turn input lines into arguments            \\
	\cmd{... | xargs -n 1 cmd}       & One argument per invocation                \\
	\cmd{... | xargs -I\{\} cp \{\} /bak} & Place the argument somewhere specific \\
	\cmd{... | xargs -P 4 cmd}       & Run four invocations in parallel           \\
	\cmd{find . -print0 | xargs -0}  & Separate with NUL bytes, not whitespace    \\
\end{keys}

\begin{note}[Filenames with spaces]
	By default \cmd{xargs} splits on whitespace, so a file named
	\file{my file.txt} arrives as two arguments and the wrong things get deleted. The fix is the
	\cmd{-print0}/\cmd{-0} pair: NUL is the one byte a filename cannot contain.
	Make it a habit --- \cmd{find . -name '*.log' -print0 | xargs -0 rm} --- or
	sidestep the issue entirely with \cmd{find -exec}
	(\autoref{sec:sh-find}).
\end{note}

\section{Here-Documents and Here-Strings}\label{sec:sh-heredoc}

When the input you want to pipe in is text you are writing right now, there is
no need for a file.

\begin{bashcode}
# here-document: everything up to the terminator becomes stdin
cat > config.yml <<EOF
name: $PROJECT      # expanded, because EOF is unquoted
port: 8080
EOF

# quoted terminator: nothing is expanded
cat > script.sh <<'EOF'
echo "$HOME stays literal"
EOF

# <<- strips leading TABS, so the heredoc can be indented with its block
if true; then
    cat <<-EOF
    indented in the source, flush left in the output
    EOF
fi

# here-string: one line, no terminator needed
grep -q root <<< "$line"
\end{bashcode}

\begin{moral}
	Quote the terminator --- \texttt{<{}<'EOF'} --- whenever the body contains
	\texttt{\$} signs you want to survive, which is every time you are generating a
	shell script, a Makefile, or anything else with variables of its own.
\end{moral}

\section{Command and Process Substitution}\label{sec:sh-subst}

\begin{keys}[0.26]
	\texttt{\$(cmd)}   & \vocab{Command substitution}: replaced by \texttt{cmd}'s output \\
	\texttt{`cmd`}     & The old spelling. Does not nest, and escaping inside it is unreadable --- do not use it \\
	\texttt{<(cmd)}    & \vocab{Process substitution}: replaced by a \emph{filename} that yields \texttt{cmd}'s output \\
	\texttt{>(cmd)}    & The same, for writing into                                      \\
	\texttt{\$((expr))} & \vocab{Arithmetic expansion}                                   \\
\end{keys}

\begin{shell}
$ echo "built on $(date +%F) by $(whoami)"
built on 2026-09-14 by konyi
$ files=$(ls *.tex | wc -l)
$ echo $((files * 2))
16
\end{shell}

Process substitution is the less known of the two, and it solves a real problem:
comparing the output of two commands, when \cmd{diff} insists on filenames.

\begin{shell}
$ diff <(sort a.txt) <(sort b.txt)
$ diff <(ls dir1) <(ls dir2)
\end{shell}

\begin{note}
	Without it you would write two temporary files and remember to delete them.
	\texttt{<(cmd)} expands to something like \file{/dev/fd/63} --- a real path,
	which \cmd{diff} opens and reads perfectly happily. It is a Bash and Zsh
	feature, not POSIX, so a script declaring \cmd{\#!/bin/sh} cannot use it.
\end{note}

\begin{moral}
	Always quote command substitution: \texttt{"\$(cmd)"}. Unquoted, its output
	goes through word splitting and globbing like any other expansion
	(\autoref{sec:sh-anatomy}), and a filename with a space in it will take your
	script apart.
\end{moral}

\section{Exit Status and the Logical Operators}\label{sec:sh-status}

Every command returns a number when it finishes. Zero means success; anything
else is a failure, and which number often says which failure.

\begin{keys}[0.20]
	\texttt{\$?}          & The exit status of the last command                     \\
	\texttt{0}            & Success --- the only value that means it                \\
	\texttt{1}            & General failure                                         \\
	\texttt{2}            & Usually a usage error                                   \\
	\texttt{126}          & Found, but not executable                               \\
	\texttt{127}          & Command not found                                       \\
	\texttt{130}          & Terminated by \key{<C-c>} --- that is $128 + 2$, the signal number \\
\end{keys}

\begin{note}
	Zero-means-success looks backwards until you see why: there is one way to
	succeed and many ways to fail, so the failures get the numbers.
\end{note}

Statuses are what the logical operators test:

\begin{keys}[0.26]
	\texttt{a ; b}     & Run \texttt{a}, then \texttt{b}, regardless                  \\
	\texttt{a \&\& b}  & Run \texttt{b} only if \texttt{a} \emph{succeeded}           \\
	\texttt{a || b}    & Run \texttt{b} only if \texttt{a} \emph{failed}              \\
	\texttt{!\ a}      & Invert the status                                            \\
	\texttt{a \&}      & Run \texttt{a} in the background (\autoref{sec:sh-jobs})     \\
\end{keys}

\begin{eg}
\begin{shell}
$ mkdir build && cd build           # only cd if the mkdir worked
$ make || echo "build failed"       # report only on failure
$ grep -q TODO file && echo "still unfinished"
$ command -v rg >/dev/null || echo "ripgrep is not installed"
\end{shell}
	The last is the standard way for a script to check whether a tool exists
	before relying on it.
\end{eg}

\begin{moral}
	\texttt{\&\&} between commands is the cheapest error handling there is. In a
	script, \cmd{cd build \&\& make} cannot run \cmd{make} in the wrong directory,
	whereas \cmd{cd build; make} happily will --- and that difference has destroyed
	real work.
\end{moral}
